%%%%%%%%%%  scan idx 369  |  folio 356  %%%%%%%%%%
et on va
montrer 6.10 d'abord pour \(X_0\) et \(N'_0=N_0+d\) (qui vérifient bien les hypothèses de 6.10). Compte tenu de la remarque faite plus haut sur la filtration de \(R_0\), il suffit pour cela de montrer que :
\[
  \gamma^{d+i}(N_0,X_0)\gamma^d(N'_0)
  -
  (-1)^{d+i-1}(d+i-1)!X_0\gamma^d(N'_0)
  \in\Fil^{d+i+1}R_0 .
\]
Or on a \(\lambda_{-1}(N_0)=(-1)^d\gamma^d(N'_0)\)~; cette relation s'écrit donc :
\[
  (-1)^d\gamma^{d+i}(X_0\ (-1)^d\gamma^d(N'_0))
  -
  (-1)^{d+i-1}(d+i-1)!X_0\gamma^d(N'_0)
  \in\Fil^{d+i+1}R_0 .
\]
Posons \(z=(-1)^dX_0\gamma^d(N'_0)\)~; alors \(z\in\Fil^{d+i}R_0\), et on est ramené à montrer :
\[
  \gamma^{d+i}(z)-(-1)^{d+i-1}(d+i-1)!z\in\Fil^{d+i+1}R_0 .
\]
Or ceci résulte de 6.9.

Pour déduire 6.10 de ce cas particulier, on considère le \(K\)-\(\lambda\)-homomorphisme \(R_0\to R\) qui envoie \(N'_0\) sur \(N-d\) et \(X_{\alpha,j}\) sur \(x_{\alpha,j}\) pour tout \(\alpha\) et tout \(j\) (4.7). Cet homomorphisme est compatible aux \(\lambda\)-filtrations, et envoie \(N_0\) sur \(N\), et \(X_0\) sur \(x\), d'où le résultat.

\textbf{Proposition 6.11.} On suppose que \(K\) est de caractéristique nulle. Soit \(n\ge0\) un entier. Alors, le foncteur \(\Ch\) est une équivalence de catégories de la catégorie des \(K\)-algèbres graduées \(A\) telles que \(A_0=K\), \(A_k=0\) pour \(k>n\), dans la catégorie des \(K\)-\(\lambda\)-algèbres augmentées \(\Lambda\) telles que \(\Fil^{n+1}\Lambda=0\)~; un foncteur quasi-inverse est donné par \(\Lambda\rightsquigarrow\Gr\Lambda\).

Soit \(A\) une \(K\)-algèbre graduée telle que \(A_k=0\) pour \(k>n\), munie de la filtration associée à sa graduation~; alors \((\Ch A)_{n+1}=0\), donc \(\Fil^{n+1}\Ch A=0\) (6.6.5). Réciproquement, si \(\Lambda\) est une \(K\)-\(\lambda\)-algèbre augmentée telle que \(\Fil^{n+1}\Lambda=0\), alors \(\Gr^0\Lambda=K\), \(\Gr^k\Lambda=0\) pour \(k>n\).

On va définir des isomorphismes de foncteurs :
\[
  A\longrightarrow \Gr(\Ch A),
  \qquad
  \Lambda\longrightarrow \Ch(\Gr\Lambda).
\]
